% >>> typage-template fduthesis@2026.10.1 variant=doctor
\documentclass[type=doctor]{fduthesis}
% <<< typage-template
\fdusetup{
  style = {
    font         = times,
    cjk-font     = fandol,
    bib-backend  = bibtex,
    bib-resource = {refs.bib},
  },
  info = {
    title       = {论文题目},
    title*      = {Thesis Title},
    author      = {作者姓名},
    author*     = {Author Name},
    supervisor  = {导师姓名\quad 教授},
    major       = {某某专业},
    degree      = academic,
    department  = {某某学院},
    student-id  = {00000000000},
    keywords    = {关键词一, 关键词二, 关键词三},
    keywords*   = {keyword one, keyword two, keyword three},
  },
}
\usepackage{booktabs}

\begin{document}

\frontmatter
\tableofcontents

\begin{abstract}
  用一段话概述研究的问题、方法与主要结论。摘要应当独立成篇，
不使用引用和未定义的缩写。
\end{abstract}

\begin{abstract*}
  Summarize the problem, the method and the main results in one
paragraph. The abstract stands on its own: no citations and no undefined
abbreviations.
\end{abstract*}

\mainmatter

\chapter{引言}
\label{chap:intro}

\section{研究背景}
介绍研究问题的背景与意义，并综述已有工作。引用文献时使用
\cite{knuth1984texbook,lamport1994latex}。

\section{本文的主要工作}
列出本文的主要贡献，并说明全文的结构。

\chapter{研究方法}
\label{chap:method}

描述所用方法。需要引用的公式应当编号，例如
\begin{equation}
  \label{eq:objective}
  \min_{\theta} \; \frac{1}{n} \sum_{i=1}^{n} \ell\bigl(f_\theta(x_i), y_i\bigr) .
\end{equation}

\chapter{实验与分析}
\label{chap:experiments}

表~\ref{tab:results} 给出了求解式~\eqref{eq:objective} 的结果。

\begin{table}[!htb]
  \centering
  \caption{示例数据集上的结果（五次运行的均值 $\pm$ 标准差）}
  \label{tab:results}
  \begin{tabular}{lcc}
    \toprule
    方法     & 准确率（\%）     & 时间（s） \\
    \midrule
    基线方法 & 81.2 $\pm$ 0.4 & 12.4      \\
    本文方法 & 84.7 $\pm$ 0.3 & 10.9      \\
    \bottomrule
  \end{tabular}
\end{table}

\chapter{总结与展望}
总结全文的工作，并指出今后的研究方向。

\appendix

\chapter{补充材料}
附录用于放置正文以外的推导、数据与程序说明。

\backmatter

\printbibliography

\begin{acknowledgements}
  感谢导师的指导，以及在研究和论文写作中给予帮助的老师与同学。
\end{acknowledgements}

\end{document}
